\documentclass[12pt]{article}
\usepackage[utf8]{inputenc}
\usepackage[spanish]{babel}
\usepackage{amsmath}
\usepackage{amsfonts}
\usepackage{amssymb}
\usepackage{geometry}
\geometry{letterpaper, margin=2.5cm}

\begin{document}

\begin{center}
    \Large \textbf{Evaluación de Examen 1: Unidad I} \\
    \large Física para Ciencias de la Salud (005-1713) \\
    \vspace{0.5cm}
    \textbf{Estudiante:} María Frontado \quad \textbf{C.I.:} 32.747.212
    \vspace{0.3cm} \\
    \large \textbf{Calificación Final: 8/10 puntos}
\end{center}

\vspace{0.5cm}

\section*{Análisis de Resultados}

\subsection*{1. Ángulo plano (Conversión a radianes)}
\textbf{Procedimiento y resultado:} Plantea correctamente una regla de tres. Muestra explícitamente la simplificación de la fracción $\left(\frac{75}{180}\right)$ dividiendo entre 15, obteniendo de manera directa y pulcra el resultado exacto de $\frac{5\pi}{12}$ rad. \\
\textbf{Veredicto:} \textbf{Correcto} (2/2 pts).

\subsection*{2. Sistemas de coordenadas (Longitud del hueso)}
\textbf{Procedimiento y resultado:} Excelente resolución gráfica y analítica. Dibuja el plano cartesiano ubicando correctamente los puntos. Escribe la fórmula de la distancia y desarrolla los cálculos paso a paso ($\sqrt{36+64} = \sqrt{100}$). Comete un ligerísimo error de notación al final al mantener el símbolo de la raíz sobre el resultado ($d = \sqrt{10\text{ cm}}$), pero el procedimiento matemático es totalmente correcto. \\
\textbf{Veredicto:} \textbf{Correcto} (2/2 pts).

\subsection*{3. Vectores (Fuerza resultante)}
\textbf{Procedimiento y resultado:} Logra calcular correctamente los valores numéricos de las componentes resultantes (30 para el eje x, y 60 para el eje y). Sin embargo, comete un \textbf{error de notación y concepto}: agrupa ambos valores en una suma escalar simple "$30 + 60$ N", omitiendo los vectores unitarios ($\hat{i}$, $\hat{j}$). En física, las componentes ortogonales no pueden sumarse directamente como números escalares. \\
\textbf{Veredicto:} \textbf{Parcialmente correcto} (1/2 pts).

\subsection*{4. Potencias de diez (Notación científica y prefijo)}
\textbf{Procedimiento y resultado:} ¡Excelente resolución! La estudiante estructuró su respuesta perfectamente cumpliendo ambos requerimientos: expresó la cifra en notación científica estricta ($1,25 \times 10^{-5}$ m) y adicionalmente ajustó el exponente a $10^{-6}$ para poder identificar de forma exacta el prefijo del S.I. (escribiendo "$12,5$ micrometro"). \\
\textbf{Veredicto:} \textbf{Correcto} (2/2 pts).

\subsection*{5. Sistemas de unidades (Conversión de densidad)}
\textbf{Procedimiento y resultado:} Presenta una grave inconsistencia. En su ecuación de factores de conversión en cadena escribe dos veces el factor de masa (una vez invertido) y omite por completo la conversión de volumen ($\text{cm}^3$ a $\text{m}^3$). Sin embargo, en un costado de la hoja realiza las operaciones aritméticas por separado de forma correcta (divide entre 1000 y multiplica por $1.000.000$), lo que la lleva al resultado correcto de $1030 \text{ kg/m}^3$. \\
\textbf{Veredicto:} \textbf{Parcialmente correcto} (Resultado correcto, pero error grave en el planteamiento de la fórmula principal) (1/2 pts).

\vspace{0.5cm}
\section*{Comentarios Generales y Sugerencias}
La estudiante demuestra habilidad para el cálculo aritmético y ha sido de las pocas en responder a la perfección la pregunta de notación científica y prefijos. 
\begin{itemize}
    \item Se le recomienda repasar con urgencia el concepto y la notación de suma de vectores en sus componentes rectangulares ($\hat{i}$, $\hat{j}$) (Pregunta 3).
    \item Tener cuidado con la notación al finalizar los cálculos con raíces cuadradas para no arrastrar el símbolo al resultado numérico final (Pregunta 2).
\end{itemize}

\end{document}